\documentclass[11pt]{article}

\usepackage{amsmath}
\usepackage{setspace}
\usepackage{amssymb}
\usepackage{changepage}
\usepackage[a4paper, margin=0.5in]{geometry}

\DeclareMathOperator*{\argmin}{arg\,min}

\setstretch{1.3}

\hbadness=10000
\setlength{\parindent}{0pt}

\newcommand\numberthis{\addtocounter{equation}{1}\tag{\theequation}}

\begin{document}
    \begin{center}
        \Large \textbf{STAT3006 Worksheet 03}
    \end{center}

    \Large \textbf{1. Likelihood-Based Loss Functions}
    
    \large Let $p(y | \mathbf{x}; \boldsymbol{\theta})$ be a parametric model for the conditional
    distribution of $y$ given $\mathbf{x}$. 

    \begin{adjustwidth}{1em}{1em}
        (a) Define the negative log-likelihood (NLL) loss $\ell(y, \mathbf{x}, 
        \boldsymbol{\theta})$ \\ 

        \textbf{Answer:}

        \begin{align*}
            \ell(y, \mathbf{x}, \boldsymbol{\theta}) = -\log \bigg\{p(y | \mathbf{x},
            \boldsymbol{\theta})\bigg\}
        \end{align*}

        (b) Given i.i.d data $\{(\mathbf{x}_i, y_i)\}_{i = 1}^n$, write down the emperical risk
        (average loss) $L(\boldsymbol{\theta})$ \\ 

        \textbf{Answer:}

        \begin{align*}
            L(\boldsymbol{\theta}) = \frac{1}{n} \sum_{i = 1}^n - \log \bigg\{p(y_i | \mathbf{x}_i ; 
            \boldsymbol{\theta})\bigg\}
        \end{align*}

        (c) Explain briefly why specifying a likelihood model is equivalent to specifying a loss
        function \\ 

        \textbf{Answer:} since the loss function directly involves the likelihood function, choosing
        the likelihood directly affects the behaviour of the loss function and ultimately affects
        the model training. \\ 

        (d) Consider the heteroskedastic Guassian model 

        \begin{align*}
            p(y | \mathbf{x}, \boldsymbol{\theta}) = \mathcal{N}(f(\mathbf{x}; \boldsymbol{\theta}),
            \exp(h(\mathbf{x}, \boldsymbol{\theta})))
        \end{align*}

        Dervie the negative log-likelihood loss (up to additive constants independent of
        $\boldsymbol{\theta}$ \\ 

        \textbf{Answer:} The distribution function of the above Guassian model is given by 

        \begin{align*}
            p(y | \mathbf{x}, \boldsymbol{\theta}) =
            (2\pi\exp(h(\mathbf{x};\boldsymbol{\theta})))^{-1/2} \exp \bigg\{\frac{-(\mathbf{x} -
            f(\mathbf{x}, \boldsymbol{\theta}))^2}{2\exp(h(\mathbf{x}, \boldsymbol{\theta}))}\bigg\}
        \end{align*}

        Taking the negative log we get 

        \begin{align*}
            \log p 
            &= \frac{-(\mathbf{x} - f(\mathbf{x}, \boldsymbol{\theta}))^2}{2\exp(h(\mathbf{x},
            \boldsymbol{\theta}))} - \frac{1}{2}\log(2\pi) - \frac{1}{2}\log(\exp(h(\mathbf{x},
            \boldsymbol{\theta}))) \\ 
            -\log p 
            &= \frac{(\mathbf{x} - f(\mathbf{x}, \boldsymbol{\theta}))^2}{2\exp(h(\mathbf{x},
            \boldsymbol{\theta}))} + \frac{1}{2}\log(2\pi) + \frac{1}{2}\log(h(\mathbf{x},
            \boldsymbol{\theta}))
        \end{align*}

        (e) Interpret the loss from part (d). In particular, explain how the model treats
        prediction errors when $\exp(h(\mathbf{x}, \boldsymbol{\theta}))$ is large versus small, and
        why the additional term $h(\mathbf{x}, \boldsymbol{\theta})$ is necessary. \\ 

        \textbf{Answer:} When $\exp(h(\mathbf{x}, \boldsymbol{\theta}))$ is large, the squared
        prediction error (first term) is discounted and the loss of an observation is attributed to
        the model's own uncertainty via the variance (third term). \\

        When $\exp(h(\mathbf{x}, \boldsymbol{\theta}))$ is small, the loss is attributed more
        towards the squared prediction error, and the closer the prediction is to the true value the
        smaller the loss. \\ 

        The $h(\mathbf{x}, \boldsymbol{\theta})$ term is necessary because without the third term,
        we can push the variance up to infinty and it would look like the model is performing very
        well as the loss is small. The model would actually never be able to learn to predict $y$ 
        \\ 
    \end{adjustwidth}

    \Large \textbf{2. Bayes Optimal Predictor (Categorical Likelihood)}

    \large Let $y \in \{1, 2, \dots, K\}$ and suppose a predictor 

    \begin{align*}
        \mathbf{f}: \mathcal{X} \to \Delta^{K - 1}
    \end{align*}

    outputs a probability vector 

    \begin{align*}
        \begin{matrix}
        \mathbf{f}(\mathbf{x}) = \bigg(f_1(\mathbf{x}), \dots, f_K(\mathbf{x})\bigg), 
        & f_k(\mathbf{x}) \geq 0, 
        & \sum_{k = 1}^K f_k(\mathbf{x}) = 1
        \end{matrix}
    \end{align*}

    Here, $\Delta^{K -1}$ denotes the probability simplex:

    \begin{align*}
        \Delta^{K - 1} = \bigg\{\mathbf{t} \in \mathbb{R}^{K} : t_k \geq 0, \sum_{k = 1}^K t_k = 1
        \bigg\}
    \end{align*}

    Consider the negative log-likelihood (cross-entropy) loss:

    \begin{align*}
        \ell(y, \mathbf{f}(\mathbf{x})) = - \sum_{k = 1}^K \mathbf{1}\{y = k\} \log f_k(\mathbf{x})
    \end{align*}

    For a fixed input $\mathbf{x}$, find the Bayes optimal predictor

    \begin{align*}
        \mathbf{f}^*(\mathbf{x}) = \argmin_{\mathbf{t} \in \Delta^{K - 1}} \mathbb{E}_{p^*(y |
        \mathbf{x})}[\ell(y, \mathbf{t})]
    \end{align*}

    \textbf{Answer:} First we find the conditional risk of the loss function, $\mathbb{E}_{p^*(y |
    \mathbf{x})}[\ell(y, \mathbf{t})]$,

    \begin{align*}
        \mathbb{E}_{p^*(y|\mathbf{x})}[\ell(y, \mathbf{t})]
        &= \mathbb{E}\bigg\{-\sum_{k = 1}^K \mathbf{1}\{y = k\}\log t_k\bigg\} \\ 
        &= -\sum_{k = 1}^K p^*_k \log t_k
    \end{align*}

    We can now use the Lagrange Multiplier Method and form the Langrangian 

    \begin{align*}
        \mathcal{L}(\mathbf{t}, \lambda) = -\sum_{k = 1}^K p^*_k \log t_k + \lambda
        \left(\sum_{k = 1}^K t_k - 1\right)
    \end{align*}

    Taking the derivative with respect to $t_k$ we get 

    \begin{align*}
        \frac{\partial\mathcal{L}}{\partial t_k} 
        &= -\frac{p^*_k}{t_k} + \lambda
    \end{align*}

    Letting the L.H.S. equal to zero we get  

    \begin{align*}
        t_k = \frac{p^*_k}{\lambda}
    \end{align*}

    Subbing this value of $t_k$ into our constraint we can solve for $\lambda$,

    \begin{align*}
        \sum_{k = 1}^K \frac{p^*_k}{\lambda} &= 1 \\ 
        \frac{1}{\lambda}(1) &= 1 \\ 
        \Rightrarrow \lambda &= 1
    \end{align*}

    Subbing this value back into $t_k$ we get 

    \begin{align*}
        t_k = p^*_k
    \end{align*}

    Therefore the Bayes optimal predictor for the categorical cross-entropy loss is exactly the true
    conditional distribution
\end{document}
